\section*{September 7, 2026: Recorded density exclusions}
\textit{Codex draft. The exclusion decision was confirmed by Jacob.}

Jacob approved excluding the Fletcher two-home parent, Eastgate 29-home parent,
and Park Place seven-home parent from density estimation because the counted
homes cannot be assigned a defensible land denominator. Their construction
records remain. The three decisions are recorded by project identifier, reason,
evidence and date in a committed source ledger. The restored exporter reads
that ledger and disables both density eligibility flags before calculating
floor-area ratio and dwelling units per acre. No land allocation is imputed.

All three observations enter both current all-construction density regressions,
which use 3,692 projects. None enters the multifamily-only regressions, which
use 822 projects. This was checked against the observations used by the fitted
models. A controlled replay of the restored exporter retained all 8,648 records
and changed only the two eligibility flags and two density measures for the three
ledger identifiers. The frozen paper input and figures remain unchanged until
the restored production chain and resulting estimates are verified.

The residential manual-review bundle producer was also restored. Its 184 rows
and 76 columns match the preserved evidence bundle byte for byte. Its inputs
include parcel histories, permit chains, building evidence and previously
recorded address matches. This verifies one stage on preserved inputs; it does
not resolve the remaining acquisition, historical zoning or source-vintage work.

Further restoration reproduced the City-building evidence for all 184 reviewed
projects and the eight accepted predecessor-parcel geometries. The condominium
review also reproduces 52 cohort-evidence rows, 57 decisions and 53 source-project
dispositions. The restored script had calculated the source dispositions but
omitted the write needed downstream; that omission is now fixed without changing
any decision. Deleting that generated output in a temporary parallel Make build
caused it to be regenerated, and the next unchanged build ran no producer.

For the preserved review scope, the replacement Assessor source returns the same
4,627 history rows. Three fields differ, all in 2026 records: one reported
construction year changes from 1895 to 2019, and two parcel-proration rates change
from one each to 0.6 and 0.4. Replaying the construction-episode evidence adds one
historical episode row, from 596 to 597, but leaves the selected evidence for all
123 reviewed projects unchanged. This comparison narrows the vintage concern
for this fixed scope; it does not establish that the changed discovery universe
or other source uses are reconciled.

Holding the current project data, alderman scores and regression specification
fixed, excluding the three ledger observations changes the estimated FAR
reduction immediately across the boundary from 12.31\% to 12.59\%. The estimated
reduction in dwelling units per acre changes from 16.01\% to 16.28\%.
The corresponding all-construction sample falls from 3,692 to 3,689.
The multifamily estimates remain exactly unchanged. This sensitivity calculation
isolates the approved exclusions; it does not predict the combined effect of
all unfinished reconstruction changes. The release audit's
\path{density_exclusion_comparison.csv} records all distance-bin estimates,
standard errors, p-values and sample sizes for both scenarios.

\paragraph{Wider review and the Lenox decision.}
The broader screen covers all 248 retained multicard projects. It finds 27
parent polygons containing current parcel points associated with other retained
projects, including 11 parents within 500 feet. These are review leads, not 27
confirmed duplicate projects. Nine projects lack a usable parcel polygon in the
preserved query file and remain unassessed by this check.

The Lenox site illustrates why a date conflict needs direct review. Its two-house
parent reports 2008, while two individually retained successors report 2019.
The original 2007 permits were reinstated in 2018, so the later date cannot be
assumed to be a reporting error. Jacob approved excluding all three records
from density until completion timing is resolved. The exporter preserves the
records and changes only their four density eligibility/outcome fields.
Adding this decision reduces the all-construction regression sample to 3,688;
the headline reductions remain 12.59\% for FAR and 16.28\% for units per acre
at displayed precision. The multifamily results remain unchanged. The audit
now reports separate scenarios for the original land decisions and all six
exclusions, with current scores and construction records held fixed.

The wider review also exposed a mismatch between a one-year window in the
project-overlap screen and a two-year window in the subsequent card matcher.
One five-townhouse project at 33rd/Prairie is represented by a 2019 parent and
five 2021 successor rows. Its permit identifies five units, and parent and successor building-area totals
differ by 0.4\%. Jacob approved a consistent two-year automatic window. The
first stage comparison adds 62 overlap edges and 15 card matches, while eight
parents change review category. Three newly required manual decisions and the
existing decisions affected by changed group membership still need reconciliation.
The final adjudicator was checked and stops with the affected project identifiers.
The component check also precedes the final single-family unit correction:
it sees 20 units in the five Prairie successors, which later become five units.
Across the expanded graph, five records in two groups differ from their later
unit classification. These comparisons do not yet establish revised density results.
Rockwell's 2006 house permits have 2005 application dates, which the archived
verification filter dropped. Jacob approved retaining all dates in the pinned
permit source. The revised stage adds 5,228 permits, including both Rockwell
house permits; all 828,237 previous records are unchanged. Existing coordinate
and processing-time validity rules remain. This restores evidence without
settling the houses' completion years or changing analysis-specific date limits.

\paragraph{Further restoration checks.}
Five more preserved outputs reproduce exactly: project-overlap evidence (313
rows), remaining residential cases (7), condominium requests (75), current parcel
links (1,789), and shared-site location review (2,699). The construction-year
zoning replay reproduces 9,609 rows after moving its embedded judgments into
explicit inputs. The reconstructed map preserves all 11,294 geometries; its
current parser changes two group fields on one RT4A polygon relative to the
preserved map. Replaying the next two zoning stages changes no project output
cells on the preserved-input comparison. Upstream ordinance extraction and
matching are still unfinished.

Ordinary CPI, parcel-sales and residential-characteristics builds now verify and
restore preserved source bytes. Separate refresh recipes write candidate files.
Their 1,339, 902,830 and 7,489,601 rows respectively reproduce the prior inputs
exactly, and their declared keys pass validation. Missing-output and failed-source
checks use a disposable fixture. Census acquisition, licensed rent-source vintage
protection, and other missing construction producers still require work. These
checks do not complete items 1--5 or authorize starting the full construction
and paper rebuilds.

\noindent\begin{minipage}{\linewidth}
\textbf{Reproduction note.}
Land decisions are committed as \texttt{79d35a3}; Lenox exclusions as
\texttt{4e62bf8}. The construction cleaning task reads
\path{adjudication/density_denominator_decisions.csv}.
Separate exporter comparisons apply the three land exclusions and the three
Lenox exclusions with other inputs fixed. Stage hashes and comparison limits
are in the release audit's \path{restored_stage_checks.csv}.
\end{minipage}
